\documentclass[11pt]{article}
\usepackage{palatino}
\usepackage{latexsym}
\usepackage{verbatim}
\usepackage{graphicx}
\usepackage{subcaption}
\usepackage{alltt}
\usepackage{url}
\usepackage[margin=1.0in]{geometry}
\usepackage{enumitem}
\usepackage{listings}
\usepackage{amsmath,proof,amsthm,amssymb}

\usepackage{hyperref}
\hypersetup{
    colorlinks=true,
    linkcolor=blue,
    filecolor=magenta,
    urlcolor=cyan,
}

\newcommand{\question}[3]
  {\bigskip \noindent
  {\bf Question #1}, {#2}, {\it(#3 points).}}

\newcommand{\extracreditquestion}[2]
  {\bigskip \noindent
   {\bf Extra Credit Question}, {#1}, {\it(#2 points).}}

\newcommand{\optionalquestion}[2]
  {\bigskip \noindent
  {\bf Question #1}, {#2}, {\it(optional).}}

\newcommand{\definition}[2]
  {\bigskip
   \begin{tabular}{p{1.5in}p{4.0in}}
        \textbf{#1} & #2 \\
        \end{tabular}
  }

\def\prop{\textsf{\,\,prop}}
\def\thm{\textsf{\,\,thm}}
\def\implies{\Rightarrow}

\title{Homework 7: Verification, SMT, and Synthesis}
\author{17-355/17-665/17-819: Program Analysis\\
        Fraser Brown and Yiliang Liang}
\date{Due: Thursday, November 19, 2026 11:59 pm EDT \\
\emph{(except for Question 3 -- see instructions)} \\[1em]
100 points \\
\emph{(+ 3\% course grade bonus for Question 3)}
}

\begin{document}

\newtheorem{theorem}{Theorem}
\newtheorem{lemma}[theorem]{Lemma}

\newcommand{\pfalse}{\mbox{false}}
\newcommand{\ptrue}{\mbox{true}}
\newcommand{\pand}{\mbox{and}}
\newcommand{\por}{\mbox{or}}
\newcommand{\pskip}{\mbox{skip}}
\newcommand{\pif}{\mbox{if}}
\newcommand{\pthen}{\mbox{then}}
\newcommand{\pelse}{\mbox{else}}
\newcommand{\pdo}{\mbox{do}}
\newcommand{\pwhile}{\mbox{while}}


\maketitle

\vspace{1ex}\noindent\textbf{Assignment Objectives:}

\begin{itemize}[labelwidth=0.7em, labelsep=0.6em, topsep=0ex, itemsep=0ex,
  parsep=0ex, leftmargin=1.0em]
\item Reason about satisfiability and EUF theory and implement a solution in SMT-LIB format.
\item Demonstrate understanding of specifications for verifying programs.
\item Develop an SMT formulation and implement the solution for an Inductive
  Synthesis problem.
\item \emph{(Optional)} Read some papers about finding bugs in the real world.
\end{itemize}

\paragraph{Handin Instructions.} This assignment has two components (Written and Coding). Questions 1 and 2a are part of the Written component, and Question 2b is part of the Coding component. The written component should be submitted on Gradescope as a file containing your answers (just like the written assignments before), and the coding component should be submitted as a GitHub repository to Gradescope. Question 3 is an optional question. See the instructions in the assignment for details.

\textbf{Note:} \emph{We will not look at your homework repository directly, but
  will only see what you have submitted to Gradescope! \textbf{Make sure that
    you (re)submit after you have completed all parts; Gradescope does not
    automatically pull new commits from GitHub.}}

\question{1}{SMT with EUF (Written)}{25}
% adapted from Q1 of https://courses.cs.washington.edu/courses/cse507/17wi/assignments/hw2.pdf

\noindent\emph{a) (10 points)} Show, using the congruence closure, whether the following Equality Logic with
Uninterpreted Functions (EUF) formula is
satisfiable. Show each step merging equivalent terms.

\[
f(g(0)) = g(f(0)) \land f(g(f(y))) = 0 \land f(y) = 0 \land g(f(0)) \neq 0
\]

\vspace{1em}
\noindent\emph{b) (15 points)} Give the SMT-LIB formula (i.e., a valid Z3 program) to prove your answer to (a) is
correct. Put your code in the answer, and show the output of Z3.

% (declare-fun f (Int) Int)
% (declare-fun g (Int) Int)
% ; (declare-fun x () Int)
% (declare-fun y () Int)
% (assert (= (f (g 0)) (g (f 0))))
% (assert (= (f (g (f y))) 0))
% (assert (= (f y) 0))
% (assert (not (= (g (f 0)) 0)))
% (check-sat)
% (get-model)    

\newpage

\question{2}{Inductive Synthesis (Written and Coding)}{75}

This task asks you to write a simple synthesizer for an expression that satisfies a
collection of input/output pairs. The instruction set consists of only four
binary operators for bitvectors: \texttt{multiply}, \texttt{add}, \texttt{left
  bit shift}, and \texttt{bitwise or}. Your expression has four inputs \texttt{A},
\texttt{B}, \texttt{C}, \texttt{D}. Your task is to discover the right
operators using the inputs such that it evaluates to the output.

\textbf{Example.} Suppose we have the input A = 2, B = 2, C = 1, D = 1, and the
output 4. Two solutions are possible, and represented with different trees ((A +
B) * (C * D)) and (A * (B * (C * D))).


More generally, we want to discover the operators and expression tree such that
the inputs evaluate to the output.

\vspace{1em}
\noindent\emph{a) (25 points)} (Written) Develop a strategy to encode the general problem with SMT.

\vspace{1em}
\noindent\emph{b) (50 points)} (Coding) 
Make a private fork of the starter repository at
\[\text{\url{https://github.com/CMU-program-analysis-fa26/homework-7}}.\] Follow the instructions in the repository and implement your solution (using your SMT
formulation) to synthesize the expression satisfying the input/output pairs
using Z3. The input/output pairs are provided in the starter repository. Inputs
correspond to the format (A, B, C, D).

You are free to use any programming language that interfaces with Z3. APIs exist
for C, Java, Python, OCaml, etc. See
\href{http://z3prover.github.io/api/html/index.html}{here}.  We provide starter
code using Python in the provided repository.

%\small
%\begin{verbatim}
%OUT:  35840 IN: (11, 6, 9, 4)
%OUT:  46080 IN: (9, 12, 10, 1)
%OUT:     12 IN: (2, 4, 0, 4)
%OUT:  52224 IN: (14, 8, 10, 3)
%OUT:   6656 IN: (1, 13, 9, 4)
%OUT:   2048 IN: (4, 2, 8, 0)
%OUT:  57344 IN: (15, 6, 12, 4)
%OUT:   5376 IN: (7, 12, 6, 0)
%OUT:      0 IN: (2, 0, 6, 0)
%OUT:    520 IN: (16, 4, 3, 1)
%\end{verbatim}
%\normalsize

\vspace{1em}
\noindent\textbf{Description.} You may assume that each input is used only
  \emph{once} in the expression, and there are four inputs. I.e., the solution is
  guaranteed to use each of the inputs A, B, C, and D, and will only use
  each of these values once.
% \item Each operator is used at most once.
All values (inputs and outputs) are 16-bits wide.
Note that inputs can occur in any order (\texttt{B << A} is a possible
  subexpression).

\vspace{1em}
\noindent\textbf{Test.}  Check your answer from your Z3 solution by making the
\texttt{test.py} file (in the same directory) pass. That is: translate your
synthesized expression into Python and add it to the test function.  If you
choose not to use Python, write a function like \texttt{test.py} that tests your
answer against all the input/output pairs, and include instructions to run your
solution in your \texttt{answers.pdf} writeup.

\vspace{1em}
\noindent\textbf{Guidelines.}  Permutations of binary trees correspond to the set of
possible expression trees that can be generated. Leaf nodes correspond to
inputs, and we know there are four leaf nodes since we have four inputs. Non-leaf
nodes correspond to operators. Can you think of a way to encode
permutations of binary trees with the operators in a way that Z3 can tell you
which combination of operators satisfy the output?

Consider using boolean variables that correspond to operators which toggle
whether an operator is used or not in an expression. If you set up the
constraints the right way, Z3 can tell you which operators should be used to
satisfy the solution (toggled ``on'') and which ones should be left out. This
way, the Z3 solution consists of flags that correspond to the right
operator for satisfying the input/output relation.

\newpage

\optionalquestion{3}{Paper Reading on Bug-Finding}

\emph{
Successful completion of this task may earn you extra credit, not to exceed an additional 3 percentage points for your course grade, at the discretion of the instructor(s). For example, if your course grade is 75\% and you receive full credit for this task, your grade increases to 78\%.
}

Throughout this class, we learned many program analysis and bug-finding techniques in a lot of theoretical details, but we never looked into how these techniques (and other methodologies derived from them) are used to find real bugs. This optional question asks you to read one or more research paper of your choosing from a list of pre-approved papers and topics, and prepare a short, 5-10 minutes class presentation summarizing the paper(s) you read.

\vspace{1em}
\noindent\textbf{Topic/Paper Selection.} Please use \href{https://docs.google.com/spreadsheets/d/1CNaxszvSpcCFQ9AK9qFv95xHyorUQ9-02bErdetcyjE/edit?usp=sharing}{this link} to select a topic\footnote{Log in using your CMU email.} by putting your name next to the topic that you choose. Each topic has one or more papers/tools associated with it. You do not need to read all of them, but you are certainly welcomed to read and discuss more than one.

If you have a topic in mind that is not listed (and some paper(s) in that topic that you would like to read), please discuss it with the instructor(s) for approval before proceeding.

\vspace{1em}
\noindent\textbf{Reading Paper(s).} Read at least one paper / article associated with your chosen topic listed in the spreadsheet. You may read more if you wish. Pay particular attention to the problem being addressed, the methodology used, and the key results.

A few questions to consider while reading the paper(s) include:
\begin{itemize}
  \item Why do previous approaches fail to address the problem effectively?
  \item What does this paper do differently to address the problem? What are the novel contributions? Why do these contributions work? \emph{For example, just identifying that a paper implements fuzzing is not enough. What specific techniques of fuzzing did they use, and why did this technique work well for their specific problem domain? What properties of the system that they are analyzing make these techniques particularly applicable? How did they guide the fuzzing process? What did these techniques not achieve? These are all more insightful than ``they used fuzzing.''}
  \item What are the key findings? What bugs did they find using these techniques? Can you illustrate some of them?
\end{itemize}

Focus on the high-level ideas without getting bogged down by low-level technical and implementation details. A few of the topics also have some associated tools. You are welcomed, but not required, to explore and try out these tools.

\vspace{1em}
\noindent\textbf{Presentation.} Summarize your findings into a short, 5-10 minute class presentation. You may (but are not required to) use slidedecks.

\vspace{1em}
\noindent\textbf{Deadlines (for Question 3).}
The deadline for this question differs from the other questions in this assignment.
\begin{itemize}
  \item Topic selection: Tuesday, Dec. 1 at 11:59pm ET.
  \item Presentation: recitation section on Friday, Dec. 4, 10-10:50am ET.
\end{itemize}

No late days allowed for topic selection and presentation.

\vspace{1em}
\noindent\textbf{Acknowledgements.} This assignment (and particularly the paper selection) is adapted from Rohan Padhye's course ``Fantastic Bugs and How to Find Them'' (17-712) which goes through much of the same papers, but in much greater depth. It's a fun class, you should take it when it's offered.

\end{document}

%%% Local Variables:
%%% mode: latex
%%% TeX-master: t
%%% End:
